\section{Introduction}
\label{sec:intro}

Antibiotic resistance is an escalating clinical problem, and biofilms (dense, surface-attached bacterial communities) carry a disproportionate share of it. Biofilms are responsible for a large fraction of chronic infections and are notoriously hard to clear with antibiotics; standard regimens that work in well-mixed culture often fail when the target population is organized as a biofilm. This has driven interest in behavior-modifying interventions, in particular compounds that induce bacterial dispersal away from biofilms toward a planktonic, or more well-mixed state. Such interventions do not kill cells; they simply change how cells are arranged in space. The hope is that disrupting biofilm spatial structure reduces the population's capacity to evolve resistance under antibiotic pressure, essentially disarming the bacteria such that they will be more easily cleared by traditional antibiotic treatments and the host immune system. Knowing whether that hope is justified requires understanding what the spatial structure of a biofilm actually does to a population's evolutionary dynamics in the first place. Does biofilm formation help the population evolve resistance, hurt it, or some combination of both?

Biofilms have a clear spatial asymmetry: cells on the outer surface are exposed to the drug, while cells in the deeper interior are shielded from it. The shielded interior is a spatial refuge, a region of the population where the antibiotic selection pressure does not reach. We ask whether that refuge makes the population more likely to undergo \emph{evolutionary rescue}, the outcome in which a declining population avoids extinction because a resistant mutant arises and takes over in time. Concretely: does a biofilm with a shielded interior rescue more often than a well-mixed population of the same initial size?

The naive prediction is that the refuge provided by biofilm structures should help facilitate evolutionary rescue. The biofilm creates a gradient of antibiotic concentration, with the shielded interior experiencing lower levels of the drug. This shielded interior lets the population persist for longer under treatment, and a longer overall population-level lifespan means more cell divisions, which means more opportunities for a resistant mutation to arise. More mutations should mean more rescue. But this prediction has at least two distinct ways it could be wrong. First, a mutation that arises in the shielded interior does nothing on its own: there is no selection pressure inside the refuge to make a resistant cell take over, so a core mutation matters only if it eventually reaches the drug-exposed surface in time to grow under selection. The refuge might generate mutations the population can never actually use. Second, the very mechanism that maintains the refuge, namely the replacement of dying surface cells by cells flowing outward from the interior, could act against a new resistant cell trying to establish on the surface by ensuring it is surrounded by a continuously refreshed wild-type majority. The same spatial structure that buys extra time may also make each individual mutation that arises less useful. Whether the biofilm is a net help or a net hindrance therefore depends on how these effects balance.

We measure that balance directly by simulating a minimum-viable biofilm (a one-dimensional chain of cells with a drug-exposed edge and a shielded core) under antibiotic pressure, and compare its rescue probability against an otherwise identical well-mixed population. By varying the antibiotic dose and the metabolic activity of the interior independently, we ask whether the biofilm's spatial structure is uniformly helpful, uniformly harmful, or, as the results below will show, contingent on both the selection pressure and on what is happening inside the refuge.
